\documentclass[12pt]{article}
\usepackage[utf8]{inputenc}
\usepackage[T1]{fontenc}
\usepackage[margin=1in]{geometry}
\usepackage{setspace}
\usepackage{xcolor}
\usepackage[colorlinks=true, urlcolor=blue!60!black, linkcolor=blue!60!black]{hyperref}
\setlength{\parindent}{0em}
\setlength{\parskip}{0.6em}
\renewcommand{\baselinestretch}{1.0}
\begin{document}
\begin{center}
{\LARGE \textit{Moody v.\ NetChoice, LLC}, 603 U.S.\ 707 (2024)\par}
\vspace{0.5em}
{\large Case Brief --- Two-Page Edition\par}
\end{center}

\section{Memory Jogger}
Facial challenges to Florida and Texas social-media laws were remanded; curating feeds is protected editorial expression (\href{https://supreme.justia.com/cases/federal/us/603/22-277/}{Justia}).

\section{Detailed Case Facts}
In 2021 Florida and Texas barred large platforms from ``censoring'' users or posts (Texas, by viewpoint; Florida, for candidates and journalists) and required explanations for moderation decisions. NetChoice, representing Facebook, YouTube, and others, brought facial challenges.

\section{Procedural History}
The Eleventh Circuit upheld an injunction against Florida's law; the Fifth Circuit upheld Texas's law, calling moderation ``not speech.''

\section{Judicial Votes}
Vacated, 9--0. Kagan wrote for Roberts, Sotomayor, Kavanaugh, Barrett, and in part Jackson; Barrett, Jackson, Thomas, and Alito (with Thomas and Gorsuch) concurred.

\section{Holding}
Neither court applied the facial standard, which requires mapping each law's full scope and weighing unconstitutional against legitimate applications. Curating third-party posts is protected expression, and rebalancing speech is not a valid interest.

\section{Analysis of Opinions}
\textbf{Majority:} The courts examined only two feeds, not direct messaging, email, or payment apps. Under \textit{Tornillo}, \textit{Turner}, and \textit{Hurley}, compiling and labeling speech is expressive; Texas's goal of correcting the ``mix of speech'' fails even intermediate scrutiny. \textbf{Barrett:} AI moderation and foreign ownership remain open. \textbf{Alito:} Part III is ``nonbinding dicta,'' and platforms may resemble common carriers.

\section{Examples: Future Applications}
\textbf{Same-Side:} a must-carry rule for candidates' posts, as applied to a platform with misinformation rules. \textbf{Opposite-Side:} barring payment apps from refusing transfers by political affiliation. \textbf{Fence-Sitter:} a feed ranked only by an engagement algorithm.

\section{Critique}
Three Justices called Part III dicta; the algorithm question is open; and the facial standard raises litigation costs.

\section{Key Quotations}
\begin{quote}
``Deciding on the third-party speech that will be included in or excluded from a compilation---and then organizing and presenting the included items---is expressive activity of its own. And that activity results in a distinctive expressive product.''
\end{quote}
\noindent --- Justice Kagan, for the Court, 603 U.S.\ at 731.

\end{document}
